\documentclass{article}
\usepackage{amsmath}
\usepackage{graphicx}
\usepackage{hyperref}
\usepackage{cite}

\title{Complex QA and Language Models Hybrid Architectures Survey}
\author{Your Name}
\date{\today}

\begin{document}

\maketitle

\begin{abstract}
This paper reviews the state-of-the-art of language models architectures and strategies for 'complex' question-answering (CQA) with a focus on hybridization. Large language models (LLMs) demonstrate efficacy in leveraging public data for standard tasks. However, addressing complex, domain-specific questions requires sophisticated architectures, domain adaptation, multi-step reasoning, and data sensitivity considerations. This survey examines recent advancements, including contributions from projects like ChatGPT, GALACTICA, BIG, BLOOM, and HELM. We specifically explore challenges in complex QA such as long-form and non-factoid QA, hallucinations, and explainability. Current solutions and research trends are analyzed, considering hybrid LLMs, neuro-symbolic systems, and human-in-the-loop methods.
\end{abstract}

\tableofcontents

\section{Introduction}
Question-answering (QA) systems have seen tremendous advancements, particularly with the advent of large language models (LLMs) such as GPT-3, BERT, and their variants. These models are adept at processing and generating human-like text, making them suitable for a wide range of applications. However, when faced with complex questions that require deep understanding, domain-specific knowledge, and multi-step reasoning, their performance can be limited.

Complex question-answering (CQA) entails answering questions that often require synthesis across various domains, reflecting human-like reasoning, and handle nuanced contexts. An example of such a question would be, "How does the concept of personal freedom vary between different cultures?" or "What is the best mix of power generation methods to reduce climate change?" Addressing these types of questions goes beyond the scope of standard QA systems and presents unique challenges.

In this survey, we explore the need for hybrid architectures that integrate LLMs with other AI techniques to enhance their capability in handling complex QA tasks. We review the foundational skills required for such systems, the methodologies employed for evaluation, and delve into specific challenges such as domain adaptation, hallucinations, and explainability. 

\section{Required Skills and Evaluation Techniques}
To effectively tackle CQA, several skills are essential:
- **Domain Adaptation**: Transferring knowledge from one domain to another.
- **Multi-Step Reasoning**: Combining several pieces of evidence to answer a question.
- **Long-Form Answer Generation**: Producing detailed and comprehensive answers.
- **Handling Sensitivity and Safety**: Ensuring data used is secure and model outputs are safe.

Evaluation techniques for CQA systems must be stringent and multi-faceted:
- **Accuracy and Precision**: Standard metrics for QA accuracy.
- **Fairness**: Ensuring the system does not perpetuate biases.
- **Robustness**: The ability to handle adversarial inputs and remain reliable.
- **Explainability**: Providing interpretable and justifiable outputs.

Research papers such as BIG, BLOOM, and HELM provide comprehensive benchmarks and analyses to evaluate LLMs' performance (Zhang et al., 2022; Scao et al., 2022; Chang et al., 2022).

\section{Challenges in Complex QA}
\subsection{Domain Adaptation}
Domain adaptation is crucial for CQA as it allows models to apply general knowledge to specific contexts. Techniques such as fine-tuning pre-trained models on domain-specific data have shown promise (Howard et al., 2018). However, challenges remain regarding the scarcity of labeled data and the transferability of knowledge across vastly different domains (Ruder et al., 2019).

\subsection{Multi-Step QA}
Effective multi-step reasoning requires decomposing questions into manageable sub-questions and synthesizing answers. Recent approaches involve neural-symbolic reasoning and program synthesis, where LLMs generate and execute programs to derive answers (Chen et al., 2020; Nye et al., 2021).

\subsection{Long-Form and Non-Factoid QA}
Long-form QA requires generating detailed responses, often synthesizing information across multiple sources. Research on generative models like T5 and PEGASUS highlights strategies for generating coherent and contextually rich answers (Raffel et al., 2020; Zhang et al., 2020). 

Non-factoid QA, involving questions that do not have a straightforward factual answer, presents additional challenges. This includes handling open-ended questions and providing human-like reasoning (Kwiatkowski et al., 2019).

\subsection{Hallucinations}
Hallucinations refer to instances where models produce outputs that are not grounded in the input data. This is a significant challenge in LLMs as they can confidently generate incorrect information. Researchers propose integrating verification mechanisms and using structured knowledge bases to mitigate this (Maynez et al., 2020).

\subsection{Safety and Multi-Sensitivity Data Protection}
Ensuring the safe use of AI in handling sensitive data and generating ethical responses is paramount. Techniques such as differential privacy and robust data handling protocols are being researched to address these issues (Abadi et al., 2016).

\section{Hybrid LLM Architectures and Trending Solutions}
Hybrid architectures combine LLMs with other AI techniques for enhanced capability. Notable approaches include:
- **Neuro-Symbolic Systems**: Integrating neural networks with symbolic reasoning to improve interpretability and reasoning (Garcez et al., 2019).
- **Human-in-the-Loop**: Using active human reinforcement to guide and correct model outputs, ensuring higher accuracy and relevance (Hancock et al., 2018).
- **Program Synthesis**: Where models generate small programs or logic to solve parts of the questions before combining the results (Nye et al., 2021).

Iterated decomposition techniques involve breaking down complex questions into simpler sub-questions that can be more easily answered (Talmor et al., 2018).

\section{Conclusion}
The field of complex QA is rapidly evolving, with significant advances driven by hybrid architectures that integrate LLMs with other AI methodologies. Addressing the multifaceted nature of complex questions requires innovations in domain adaptation, multi-step reasoning, long-form QA, and robust safety mechanisms. As research progresses, these systems are expected to become more adept at providing accurate, reliable, and nuanced answers to complex queries.

\bibliographystyle{plain}
\bibliography{prompt1_try2_4o}

\end{document}